%\SweaveUTF8



\setkeys{Gin}{width=\textwidth}

\begin{Problem}[Simple random sampling]
    We use the data set \code{sepe.dat} from Problem Set 1.

\begin{Subtask}
    %---------------------------------------------------------------------------
    \item Draw a simple random sample without replacement (\SRS) of size $n =
        200$ from the \code{sepe.dat} population, which is of size $N = 800$.
        Estimate
        \begin{enumerate}[i)]
            \item the mean sepeironmental protection expenditure per
                establishment, $\bar{y}_U$;
            \item total expenditure, $\ty$;
            \item mean expenditure per employee among establishments with at
                least one employee,
                \begin{equation*}
                    \bar{q}_{U+} = \frac{1}{N_+}\sum_{i \in U_+}\frac{y_i}{x_i},
                    \qquad
                    U_+ = \{i \in U:x_i>0\}.
                \end{equation*}
        \end{enumerate}
        Here, $y_i$ denotes expenditure and $x_i$ the number of full-time
        employees of establishment $i$. Also estimate the standard error of
        each estimator.

        Compare the estimated standard errors with the corresponding
        population-based standard errors. For the sample mean, also compare
        the estimated variance with the Monte Carlo (MC) variance from Problem
        1.3a).

        \begin{Rhint}
            Write down the three estimators before doing the calculations in R.
            To reproduce the solution, use \code{set.seed(11)} before drawing
            the sample. A sample of indices can then be obtained with

\begin{Scode}
N <- NROW(sepe)
in_sample <- sample(N, 200)
sepe_sample <- sepe[in_sample, ]
\end{Scode}
        \end{Rhint}

    \begin{Solution}
        First load the population and draw the sample.

\begin{Schunk}
\begin{Sinput}
> sepe <- read.table("https://stat.ethz.ch/Teaching/Datasets/cas-das/sepe.dat",
                     header = TRUE)
\end{Sinput}
\end{Schunk}




\begin{Schunk}
\begin{Sinput}
> N <- NROW(sepe)
> set.seed(11)
> n <- 200
> in_sample <- sample(N, n)
> sepe_sample <- sepe[in_sample, ]
> fpc <- 1 - n / N
\end{Sinput}
\end{Schunk}

        \begin{enumerate}[i)]
            %------------------------------
            \item \textbf{Mean expenditure.} The estimator and its estimated
                standard error are

\begin{Schunk}
\begin{Sinput}
> m_y <- mean(sepe_sample$expend)
> m_y
\end{Sinput}
\begin{Soutput}
[1] 82508.45
\end{Soutput}
\begin{Sinput}
> se_m_y <- sqrt(fpc * var(sepe_sample$expend) / n)
> se_m_y
\end{Sinput}
\begin{Soutput}
[1] 3964.578
\end{Soutput}
\end{Schunk}

                Using the complete finite population, the standard error is

\begin{Schunk}
\begin{Sinput}
> SE_m_y <- sqrt(fpc * var(sepe$expend) / n)
> SE_m_y
\end{Sinput}
\begin{Soutput}
[1] 4170.373
\end{Soutput}
\end{Schunk}

                The Monte Carlo variance from Problem 1.3a) was 17479711; thus,
                the MC estimate of standard error is

\begin{Schunk}
\begin{Sinput}
> sqrt(17479711)
\end{Sinput}
\begin{Soutput}
[1] 4180.874
\end{Soutput}
\end{Schunk}

                The two standard errors are very close.

            %------------------------------
            \item \textbf{Total expenditure.} Since $\tys=N\mys$, estimate,
                (theoretical) standard error, and estimated standard error are:

\begin{Schunk}
\begin{Sinput}
> t_y <- N * m_y
> t_y
\end{Sinput}
\begin{Soutput}
[1] 66006760
\end{Soutput}
\begin{Sinput}
> se_t_y <- N * se_m_y
> se_t_y
\end{Sinput}
\begin{Soutput}
[1] 3171662
\end{Soutput}
\begin{Sinput}
> SE_t_y <- N * SE_m_y
> SE_m_y
\end{Sinput}
\begin{Soutput}
[1] 4170.373
\end{Soutput}
\end{Schunk}

            %------------------------------
            \item \textbf{Mean expenditure per employee.} Restrict attention
                to establishments with at least one employee. Conditional on
                the number $n_+$ of such establishments in the sample, they
                form an \SRS from $U_+$. Thus, the usual \SRS variance formula
                applies within this subpopulation.

                Restrict data and variables to cases with at least one
                employee:

\begin{Schunk}
\begin{Sinput}
> sepe_pos <- sepe[sepe$empl > 0, ]
> sepe_sample_pos <- sepe_sample[sepe_sample$empl > 0, ]
> N_pos <- NROW(sepe_pos)
> n_pos <- NROW(sepe_sample_pos)
\end{Sinput}
\end{Schunk}

                Define the expenditure-per-employee variable for establishments
                with at least one employee:

\begin{Schunk}
\begin{Sinput}
> q <- sepe_sample_pos$expend / sepe_sample_pos$empl
> q_pop <- sepe_pos$expend / sepe_pos$empl
\end{Sinput}
\end{Schunk}

\begin{Schunk}
\begin{Sinput}
> fpc_pos <- 1 - n_pos / N_pos
\end{Sinput}
\end{Schunk}

                The estimator and its estimated standard error are

\begin{Schunk}
\begin{Sinput}
> m_q <- mean(q)
> m_q
\end{Sinput}
\begin{Soutput}
[1] 1151.237
\end{Soutput}
\begin{Sinput}
> se_m_q <- sqrt(fpc_pos * var(q) / n_pos)
> se_m_q
\end{Sinput}
\begin{Soutput}
[1] 148.2974
\end{Soutput}
\end{Schunk}

                The corresponding population-based standard error is

\begin{Schunk}
\begin{Sinput}
> SE_m_q <- sqrt(fpc_pos * var(q_pop) / n_pos)
> SE_m_q
\end{Sinput}
\begin{Soutput}
[1] 176.0603
\end{Soutput}
\end{Schunk}

        \end{enumerate}
    \end{Solution}

    %---------------------------------------------------------------------------
    \vspace{0.5em}
    \item Repeat the three estimates from part a) using the R package
        \code{survey} and the sample \code{sepe_sample}.

        \begin{Rhint}
            Attach the \code{survey} package to the search path, add the
            population size $N$ to the sample data, and define the survey
            design object \code{dn\_sepe}:

\begin{Scode}
library("survey")
sepe_sample$N <- N
dn_sepe <- svydesign(ids = ~1, fpc = ~N, data = sepe_sample)
\end{Scode}

            The argument \code{ids = ~1} specifies a design without clustering.
            Supplying \code{fpc = ~N} incorporates sampling without replacement
            from a population of size $N$.

            New variables can be added to a design object with \code{update()}.
            A design object can be restricted to a subpopulation using
            \code{subset()}.

            Use \code{svymean()} to estimate population means and
            \code{svytotal()} to estimate population totals. Their basic syntax
            is
\begin{Scode}
svymean(~variable, design)
svytotal(~variable, design)
\end{Scode}
            The first argument is an R formula specifying the variable to be
            analyzed. The tilde \code{~} is part of the formula syntax; for
            example, \code{~expend} specifies the variable \code{expend}. The
            second argument is the survey design object containing the sample
            data and the information about the sampling design.

            The functions return objects containing both the estimates and
            information about their sampling variability. The corresponding
            methods \code{coef()}, \code{SE()}, and \code{cv()} extract the
            point estimates, standard errors, and coefficients of variation,
            respectively.
        \end{Rhint}

    \begin{Solution}
        Define the survey design.

\begin{Schunk}
\begin{Sinput}
> library("survey")
> sepe_sample$N <- N
> dn_sepe <- svydesign(ids = ~1, fpc = ~N, data = sepe_sample)
\end{Sinput}
\end{Schunk}

        The mean and total expenditure are obtained directly:

\begin{Schunk}
\begin{Sinput}
> svymean(~expend, dn_sepe)
\end{Sinput}
\begin{Soutput}
        mean     SE
expend 82508 3964.6
\end{Soutput}
\begin{Sinput}
> svytotal(~expend, dn_sepe)
\end{Sinput}
\begin{Soutput}
          total      SE
expend 66006760 3171662
\end{Soutput}
\end{Schunk}

        For mean expenditure per employee, restrict the design to
        establishments with at least one employee and then define the ratio.

\begin{Schunk}
\begin{Sinput}
> dn_pos <- subset(dn_sepe, empl > 0)
> dn_pos <- update(dn_pos, q = expend / empl)
> svymean(~q, dn_pos)
\end{Sinput}
\begin{Soutput}
    mean     SE
q 1151.2 148.14
\end{Soutput}
\end{Schunk}

        The estimates agree with part a), apart from possible differences in
        the variance calculation for the domain mean.
    \end{Solution}

    %---------------------------------------------------------------------------
    \vspace{0.5em}
    \item Using the design from part b), construct a 95\,\% Wald and a
        logit-based confidence interval for the population proportion $P_{31}$
        of establishments in industry class 31, identified by \code{code = 1}.
        Use \code{svyciprop()}.

        \begin{Rhint}
            A normal-based Wald interval on the probability scale is obtained
            with

\begin{Scode}
svyciprop(~code, dn_sepe, method = "mean", level = 0.95, df = Inf)
\end{Scode}

            Setting \code{df = Inf} specifies that the confidence interval
            should use the standard normal distribution rather than a Student
            $t$-distribution based on the design degrees of freedom.

            Using \code{method = "logit"} instead gives a logit-based
            confidence interval. The logit interval respects the natural bounds
            of a proportion and can better accommodate asymmetry when the
            estimated proportion is close to 0 or 1.
        \end{Rhint}

    \begin{Solution}
        The Wald interval is

\begin{Schunk}
\begin{Sinput}
> ci_p31 <- svyciprop(~code, dn_sepe, method = "mean", level = 0.95, df = Inf)
> ci_p31
\end{Sinput}
\begin{Soutput}
            2.5% 97.5%
code 0.170 0.125 0.215
\end{Soutput}
\end{Schunk}

        For comparison, a logit-based interval is

\begin{Schunk}
\begin{Sinput}
> svyciprop(~code, dn_sepe, method = "logit", level = 0.95, df = Inf)
\end{Sinput}
\begin{Soutput}
            2.5% 97.5%
code 0.170 0.129 0.220
\end{Soutput}
\end{Schunk}

        The two intervals need not be identical because they use different
        transformations.
    \end{Solution}
\end{Subtask}
\end{Problem}
